\chapter{Self-Adjoint Speed-Ups}\label{speedup}


In this appendix, we derive the result used to speed up our implementation of \Cref{outerloop} in \Cref{maxwell results}.

Let $A$ be a self-adjoint, closed, and densely defined linear operator on the Hilbert space $\H.$  For Colbrook's algorithm, we consider the folded operator $F(z) \defeq (A-zI)^*(A-zI).$ Since $A^*=A,$ we have
$F(z) = (A-\overline{z}I)(A-zI).$
Plugging in $z = x +iy$ and expanding, we obtain
\begin{align}
    F(z) &= (A-xI + iyI)(A-xI-iyI)\\
    &=(A-xI)^2 + y^2 I.
\end{align}
% Furthermore, for the injection modulus \eqref{inject}, we have
% \begin{align}
%     \|(A-zI) u \|_H^2 &= \| (A-zI) u - iyu \|_H^2 \\
%     &= \| (A-xI)u \|_H^2 + y^2 \|u\|_H^2 + 2 \operatorname{Re}\{ (A-xI)u, -iyu )_H \} \\
%     &= \| (A-xI)u \|_H^2 + y^2 \|u\|_H^2.
% \end{align}
% Taking the infimum over unit vectors, we can write \eqref{gamma} as
% \begin{equation}
%     \gamma(z, A)^2 = \gamma(x, A)^2 + y^2. 
% \end{equation}


% Next, we consider our variational folded operator of Section \ref{Fix matt}, written as $B(z) \defeq (L-zM)'M^{-1}(L-zM) = (L-\overline{z}M)M^{-1}(L-zM).$ Expanding, we have 
% \begin{align}
%     B(z) &= LM^{-1}L - (z + \overline z) L + |z|^2 M \\
%     &= LM^{-1} L - 2xL + (x^2 + y^2) M.
% \end{align}
% Notice that 
% \begin{equation}
%     B(x) = (L-xM)M^{-1}(L-xM) = LM^{-1}L - 2xL + x^2 M.
% \end{equation}
% Thus, 
% \begin{equation}
%     B(z) = B(x) + y^2 M.
% \end{equation}

Next we consider the quadratic form 
\begin{equation}
    Q_{z, A}(u, v) \defeq ((A-zI)u, (A-zI)v)_\H,
\end{equation}
for all $u, v \in D(A).$
Substituting in $z = x+iy,$ we have
\begin{align}
    Q_{z,A}(u,v)
    &= \bigl((A-xI)u-iyu,\,(A-xI)v-iyv\bigr)_\H \\
    &= \bigl((A-xI)u,(A-xI)v\bigr)_\H
       + iy\bigl((A-xI)u,v\bigr)_\H \notag\\
    &\quad
       - iy\bigl(u,(A-xI)v\bigr)_\H
       + y^2(u,v)_\H \\
    &= \bigl((A-xI)u,(A-xI)v\bigr)_\H
       + y^2(u,v)_\H \\
    &= Q_{x,A}(u,v)+y^2(u,v)_\H.
\end{align}
Taking $v=u,$ we have that
\begin{align}
    Q_{z,A}(u,u)
    &= Q_{x,A}(u,u)+y^2\|u\|_\H^2.
\end{align}
Then, taking the infimum over unit vectors, we can write \eqref{var gamma} as
\begin{equation}
    \gamma(z, A)^2 = \gamma(x, A)^2 + y^2. 
\end{equation}


Considering an appropriate finite element space $V_h,$ we have that
\begin{align}
    Q_{z, A}(u_h, v_h) &= Q_{x, A}(u_h, v_h) + y^2 m(u_h, v_h).
\end{align}
Thus, to seek $\lambda_h(z)$ and nonzero $u_h(z) \in V_h$ to satisfy
\begin{equation} \label{zval}
    Q_{z, A}(u_h(z), v_h) = \lambda_h(z) m(u_h(z), v_h), 
\end{equation}
for all $v_h \in V_h,$
we can instead seek $\lambda_h(x)$ and nonzero $u_h(x) \in V_h$ to satisfy
\begin{align} \label{xval}
    Q_{x, A}(u_h(x), v_h) &= (\lambda_h(z) - y^2) m(u_h(z), v_h) = \lambda_h(x) m(u_h(x), v_h), 
\end{align}
for all $v_h \in V_h.$
In other words, for each distinct $x,$ we can solve $Q_{x, A}(u_h(x), v_h) = \lambda_h(x) m(u_h(x), v_h)$ and find $\lambda_h(z)$ by adding $y^2.$ The eigenfunctions carry over as well. That is, $E_h(z) = E_h(x),$ where $E_h(\cdot)$ denotes the eigenspace for $\lambda_h(\cdot).$


% Considering an appropriate finite element space $V_h,$ we define the sesquilinear forms  $Q_{z, A}(u_h, v_h) \defeq \langle B(z) u_h, v_h \rangle$ and $m(u_h, v_h) \defeq \langle M u_h, v_h \rangle.$ Notice that 
% \begin{align}
%     Q_{z, A}(u_h, v_h) &\defeq \langle B(z) u_h, v_h \rangle \\
%     &= \langle B(x)  u_h, v_h \rangle + y^2 \langle M u_h, v_h \rangle  \\
%     &= Q_{x, A}(u_h, v_h) + y^2 m(u_h, v_h).
% \end{align}
% Thus, to seek $\lambda_h(z)$ and nonzero $u_h(z) \in V_h$ to satisfy
% \begin{equation}
%     Q_{z, A}(u_h(z), v_h) = \lambda_h(z) m(u_h(z), v_h), \qquad \forall v_h \in V_h,
% \end{equation}
% we can instead seek $\lambda_h(z)$ and nonzero $u_h(x) \in V_h$ to satisfy
% \begin{equation}
%     Q_{x, A}(u_h(x), v_h) = (\lambda_h(x) - y^2) m(u_h(z), v_h) = \lambda_h(x) m(u_h(x), v_h), \qquad \forall v_h \in V_h.
% \end{equation}
% In other words, for each distinct $x,$ we can solve $Q_{x, A}(u_h(x), v_h) = \lambda_h(x) m(u_h(x), v_h)$ and find $\lambda_h(z)$ by adding $y^2.$ The eigenfunctions carry over as well. That is, $E_h(z) = E_h(x),$ where $E_h(\cdot)$ denotes the eigenspace for $\lambda_h(\cdot).$

With this, we can conclude that for $z$ values that share the same real part, Colbrook's approximate bound on the distance to the spectrum depends only on a shift by the square of the imaginary part:
\begin{equation} \label{shift}
    \Phi_h(z, A) = \sqrt{\mu_h(z)} = \sqrt{\mu_h(x) + y^2},
\end{equation}
where $\mu_h(z)$ and $\mu_h(x)$ denote the smallest eigenvalues of \eqref{zval} and \eqref{xval}, respectively.

This result carries over to the DWR error estimate \eqref{DWR} as well. In particular,
notice that the primal residual \eqref{resid} can be written as
\begin{align}
    \rho(u_h, \lambda_h(z))(v) &\defeq Q_{z, A}(u_h , v) - \lambda_h(z) m(u_h, v) \\
    &= Q_{x, A}(u_h, v) + y^2 m(u_h, v) - \lambda_h(z) m(u_h, v) \\
    &= Q_{x, A}(u_h, v) - (\lambda_h(z) - y^2) m(u_h, v) \\
    &= Q_{x, A}(u_h, v) - \lambda_h(x) m(u_h, v) \\
    &= \rho(u_h, \lambda_h(x))(v).  
\end{align}
Since the primal residual, as well as both the primal and enriched eigenspaces, depends only on $x,$ the DWR error estimate \eqref{DWR} depends only on $x.$ Thus,
\begin{align}
    \lambda(z) - \lambda_h(z) &\approx \frac{\rho(u_h(z), \lambda_h(z))(u_h^+(z) - u_h(z))}{1 - \frac{1}{2} m(u_h^+(z) - u_h(z), u_h^+(z) - u_h(z)) }\\
    &= \frac{\rho(u_h(x), \lambda_h(x))(u_h^+(x) - u_h(x))}{1 - \frac{1}{2} m(u_h^+(x) - u_h(x), u_h^+(x) - u_h(x)) }.
\end{align}

These results allow us to drastically speed up our implementation of \Cref{outerloop}. For each distinct $x,$ we perform one expensive adaptive eigensolve at one barycentre $z_K$ that satisfies $\operatorname{Re}(z_K) = x.$ We then use \eqref{shift} to compute $\Phi_h(z_K, A)$ for every other barycentre with the same real part. The DWR error estimate computed at the representative $z_K$ is also carried over to every other barycentre with the same real part. 

% barycentre $z_K$ that shares the same real part, we perform one expensive adaptive eigensolve. Then, for each `column' (every $z_K$ that shares the same real part) in the complex plane, we compute the needed $\Phi_n(z_K, A)$ using \eqref{shift}. The error estimate computed at the chosen $z_K$ for each column carries over to every other point in that column.


% \subsection{DWR Error Estimate}

% Notice that the imaginary part of $z$ does not affect the primal residual:
% \begin{align}
%     \rho_z(v) &\defeq a_z(u_h, v) - \lambda_h(z) m(u_h, v) \\
%     &= a_x(u_h, v) + y^2 m(u_h, v) - \lambda_h(z) m(u_h, v) \\
%     &= a_x(u_h, v) - (\lambda_h(z) - y^2) m(u_h, v) \\
%     &= a_x(u_h, v) - \lambda_h(x) m(u_h, v) \\
%     &= \rho_x(v).  
% \end{align}
% Since the residual only depends on $x,$ and the eigenfunctions only depend on $x,$ the DWR error estimate depends only on $x$:
% \begin{align}
%     \lambda(z) - \lambda_h(z) &\approx \frac{\rho_z(u_h^+(z) - u_h(z))}{1 - 1/2 \| u_h^+(z) - u_h(z) \|_m^2 }\\
%     &= \frac{\rho_z(u_h^+(x) - u_h(x))}{1 - 1/2 \| u_h^+(x) - u_h(x) \|_m^2 }.
% \end{align}


% \subsection{What does this mean?}

% If correct, this means we can drastically speed up our computations. For each $x,$ we do the expensive adaptive eigensolves. Then, for each `column' (every $z$ that shares $x$) in the complex plane, we compute the needed $\Phi_n(z, A)$ using the shift by $y^2.$ The error estimate computed at $x$ carries over to every other point in that column.